\documentclass[11pt,a4paper]{article}
\usepackage[margin=0.85in]{geometry}
\usepackage{booktabs}
\usepackage{hyperref}
\usepackage{amsmath}

\title{MSML/MSAI 605 -- Milestone 2 Evaluation Report\\
\large Face Verification on LFW with FaceNet}
\author{Nithin Skantha Murugan \and Sriram Vema}
\date{March 2026}

\begin{document}
\maketitle

\section{System Overview}

This report presents the Milestone~2 evaluation of our face verification pipeline.
The system takes two face images and produces a cosine similarity score using
\textbf{FaceNet (InceptionResnetV1, pretrained on VGGFace2)} embeddings (512-dim).
A decision threshold converts this score into a binary same/different verdict.

We use the LFW dataset with deterministic splits (seed=42, 70/15/15 train/val/test)
and pair generation sorted lexicographically.
Score direction: \textbf{higher cosine similarity = more likely same person}.
All runs are tracked via our JSON-based run logger (\texttt{src/tracking/run\_logger.py}),
with a master CSV at \texttt{outputs/runs/runs.csv}.

\section{Tracked Runs and Threshold Selection}

\textbf{Threshold selection rule:} We select the threshold that \textbf{maximises F1 score
on the validation split}. This rule is stated before evaluating on the test set and applied
consistently for both baseline and post-improvement runs.

\begin{table}[h]
\centering
\small
\begin{tabular}{@{}lllcc@{}}
\toprule
\textbf{Run} & \textbf{Purpose} & \textbf{Data} & \textbf{Threshold} & \textbf{F1} \\
\midrule
Run 1 & Baseline val threshold sweep     & baseline    & --    & 0.964 (best) \\
Run 2 & Baseline threshold selection     & baseline    & 0.352  & 0.964 \\
Run 3 & Baseline test evaluation         & baseline    & 0.352  & 0.974 \\
Run 4 & Post-change val sweep            & improved\_v1 & --    & 0.963 (best) \\
Run 5 & Post-change test evaluation      & improved\_v1 & 0.352  & 0.972 \\
Run 6 & Error analysis (2 slices)        & baseline    & 0.352  & -- \\
\bottomrule
\end{tabular}
\end{table}

The selected baseline threshold is $t = 0.3518$, yielding validation
F1 = 0.964, accuracy = 0.963. The same threshold was used for post-improvement evaluation.

\section{ROC Curve and Confusion Matrix}

The ROC curve is saved at \texttt{outputs/sweeps/roc\_val.png} (200 thresholds from $-0.5$ to $1.0$).
The curve shows strong separability with the TPR rising steeply before the FPR increases.

\medskip

\noindent\textbf{Confusion matrix at $t=0.3518$ (validation set, baseline):}

\begin{table}[h]
\centering
\begin{tabular}{@{}lcc@{}}
\toprule
 & \textbf{Pred Same} & \textbf{Pred Different} \\
\midrule
\textbf{Actual Same}      & TP = 439 & FN = 18 \\
\textbf{Actual Different}  & FP = 15  & TN = 428 \\
\bottomrule
\end{tabular}
\end{table}

Accuracy = 96.3\%, Precision = 96.7\%, Recall = 96.1\%.

\section{Data-Centric Improvement}

\textbf{Problem:} The original pair set has slight label imbalance (val: 457 positive vs.\ 443 negative).
Some identities appear in disproportionately many pairs.

\textbf{Changes applied:}
(1) \textbf{Invalid image filtering} (\texttt{make\_pairs.py}): Remove pairs with missing files (0 removed--all images valid).
(2) \textbf{Identity capping} (\texttt{src/data\_centric.py}): Cap each identity to $\leq$30 pair appearances (max was 16, so 0 removed).
(3) \textbf{Label rebalancing}: Downsample majority label to 1:1 ratio.
Val: $900 \to 886$ ($-14$). Test: $901 \to 866$ ($-35$).

\textbf{Before vs.\ After (test set):}
Baseline F1 = 0.974 (901 pairs). Post-improvement F1 = 0.972 (866 pairs).
The small F1 decrease ($-0.002$) is expected: rebalancing removes ``easy'' majority-class pairs,
making the evaluation more representative rather than artificially inflated.

\section{Error Analysis}

\subsection{Slice 1: False Negatives Among Rare Identities}

\textbf{Definition:} Same-identity pairs (label=1) predicted as different, where the identity has $\leq$4 unique images in the dataset.

\textbf{Count:} 18 errors out of 426 pairs in this slice (error rate: 4.2\%).

\textbf{Examples:}
(1) Anna\_Kournikova: score = 0.157, only 3 images in dataset.
(2) Bill\_Paxton: score = 0.118, only 4 images.
(3) Charles\_Bronson: score = 0.162, only 2 images.

\textbf{Hypothesis:} Identities with few images produce less stable embeddings.
Limited intra-person variation means small pose/lighting changes shift the embedding
below the threshold, causing false rejections.

\textbf{Future improvement:} Augment rare-identity images or use a model with better few-shot generalisation.

\subsection{Slice 2: False Positives Near the Decision Boundary}

\textbf{Definition:} Different-identity pairs (label=0) predicted as same, with score in $[0.352, 0.402)$.

\textbf{Count:} 8 errors out of 8 pairs in this band (error rate: 100\%).

\textbf{Examples:}
(1) Adam\_Sandler vs.\ Saeed\_Anwar: score = 0.393.
(2) Billy\_Boyd vs.\ Sid\_Caesar: score = 0.382.
(3) Bob\_Curtis vs.\ Helen\_Clark: score = 0.368.

\textbf{Hypothesis:} These false accepts involve visually similar individuals whose embeddings
cluster just above the boundary--often similar age, gender, and conditions.

\textbf{Future improvement:} Use a more discriminative model (e.g., ArcFace) or add hard-negative mining.

\section{Conclusion}

The iteration loop confirmed that FaceNet achieves strong baseline performance
(test F1 = 0.974). Data-centric rebalancing produces a fairer evaluation ($-0.002$ F1).
The two error slices reveal failures concentrate on (1) rare identities with unstable embeddings
and (2) look-alike non-matches near the boundary.
Future work should focus on embedding quality and pair diversity rather than threshold tuning.

\end{document}
